\section{Related Work}
\subsection{Mesh Parameterization}

Disregarding mesh decomposition, mesh parameterization—the process of mapping a 3D mesh to a 2D domain while minimizing various distortion metrics (e.g., isometric, conformal, equiareal)—is a fundamental operation in UV unwrapping~\cite{sawhney2017boundary,schuller2013locally,rabinovich2017scalable}. Many methods directly optimize UV coordinates based on mesh connectivity by solving linear or nonlinear systems under fixed or free boundary conditions. These include barycentric embeddings like Tutte’s embedding~\cite{tutte1963draw}, Laplacian eigen-decomposition techniques~\cite{belkin2003laplacian,mullen2008spectral,taubin1995signal}, and distortion-minimization methods~\cite{hormann2000mips,sander2001texture,su2016area,yueh2019novel,yan2005mesh,desbrun2002intrinsic,Ray2003HierarchicalLS,ben2008conformal}, such as Least Squares Conformal Maps (LSCM)~\cite{lscm}. Some adopt local/global strategies that combine per-triangle transformations (e.g., rotations) with global stitching~\cite{sorkine2007rigid,liu2008local,alexa2023rigid}. In contrast, Angle-Based Flattening (ABF) methods~\cite{sheffer2001parameterization,sheffer2000surface,sheffer2005abf++,zayer2007linear} optimize triangle angles before converting them into UVs, typically offering better angle preservation than LSCM but at higher computational cost due to nonlinear solvers. In this paper, we adopt ABF for base unfolding, integrating it with various mesh decomposition strategies.

\subsection{Mesh Segmentation for UV Unwrapping}
Mesh segmentation is a key step in UV unwrapping, aiming to divide a complex 3D mesh into simpler patches, or charts, each of which can be flattened with minimal distortion. Existing approaches can be broadly categorized into four main strategies. Top-down methods recursively subdivide the mesh, often using spectral analysis~\cite{liu2007mesh,pothen1990partitioning,Liu2007MeshSegmentation} or cuts along feature lines~\cite{zhang2005feature} and high-curvature regions~\cite{lavoue2005new}. Bottom-up aggregation~\cite{sorkine2002bounded,lscm,kalvin1996superfaces,julius2005d, pulla2001improved, zhou2004iso, yamauchi2005mesh,takahashi2011optimized,bhargava2025mesh} grows charts from small seeds (e.g., triangles), merging them based on planarity or distortion bounds to maximize chart size under quality constraints. Cut optimization~\cite{li2018optcuts, autoCuts,pietroni2009almost,gu2002geometry,Sander2003MultiChartGI, takahashi2011optimized, carr2006rectangular} focuses on identifying optimal seam networks that minimize seam length, reduce flattening distortion, or better suit specific parameterization goals. Clustering-based methods~\cite{katz2003hierarchical,roy2023neural} group mesh elements based on geometric similarity to form coherent segments. Commonly used tools such as xatlas~\cite{xatlas}, Blender~\cite{blender}, and Open3D~\cite{zhou2018open3d} primarily adopt bottom-up aggregation strategies that rely solely on local geometric properties. In contrast, PartUV integrates high-level semantic part priors from a learning-based module with novel local geometric decomposition strategies via a top-down recursive tree search.

\subsection{Learning-Based Parameterization and Segmentation}

Representing UV mappings with neural networks has attracted attention due to their differentiability and composability. Neural Surface Maps~\cite{morreale2021neural} pioneered this direction by learning mappings across collections of surfaces. Several approaches jointly optimize UV parameterization and 3D reconstruction from multi-view images using cycle-consistency and distortion-based losses~\cite{xiang2021neutex,LearninganIsometric}. Others, such as NUVO~\cite{srinivasan2024nuvo}, solve for parameterizations directly from sampled 3D points. More recent methods~\cite{zhang2024flatten,zhao2025flexpara} design neural networks to emulate physical operations, including face cutting, UV deformation, and unwrapping. However, all the aforementioned techniques typically rely on per-shape optimization, making them computationally expensive and often requiring tens of minutes per shape. A few feedforward models address specific tasks more efficiently, such as intra-category texture transfer~\cite{chen2022auv} and low-distortion patch selection~\cite{liu2023wand}.

To enhance UV chart decomposition, we explore the integration of semantic part priors. Traditional learning-based part segmentation methods~\cite{qian2022pointnext,jiang2020pointgroup,vu2022softgroup} are restricted to closed-set categories due to the limited scale of part-annotated datasets~\cite{mo2019partnet,yi2016scalable}. Recent approaches~\cite{yang2023sam3d,xu2023sampro3d,yin2024sai3d,guo2024sam,zhou2018open3d,xu2024embodiedsam,he2024pointseg} lifting priors from 2D vision to 3D models~\cite{kirillov2023segment,li2022grounded} for open-world 3D part segmentation but rely on per-shape optimization, leading to slow runtimes and noise sensitivity. In contrast, a recent method, PartField~\cite{liu2025partfield}, introduces a feedforward model that learns part-aware feature fields for fast, hierarchical 3D part decomposition. While semantic priors intuitively benefit UV unwrapping, naïve integration can yield suboptimal results. PartUV addresses this with a novel top-down strategy that interleaves semantic segmentation and geometric flattening.